\begin{contestabstract}

针对供应商声称零配件次品率不超过标称值、企业需尽量减少检测次数的问题，本文建立精确双单侧检验与三态判定模型。以标称次品率 $p_0=0.10$ 为边界，将 $95\%$ 拒收信度与 $90\%$ 接收信度转化为两个方向相反的单侧检验；总体量未知时采用二项分布，已知有限总体且无放回抽样时采用超几何修正。结果表明，固定抽取 $22$ 件时，未发现次品则接收，发现 $1$ 至 $5$ 件次品则证据不足，发现不少于 $6$ 件次品则拒收；临界相邻概率与有限总体检验支持该判据。

针对生产过程决策，本文以两类零配件真实质量的联合后验分布为信念状态，将购件、检测、装配与拆解建模为随机最短路型马尔可夫决策过程（Markov Decision Process，MDP），并由 Bellman 方程求解状态相关策略。六种情形的最优期望利润依次为 $18.9222$、$12.0000$、$16.6533$、$14.7500$、$15.4500$、$21.6787$ 元；相较七变量固定决策基准，该策略在情形 3 将期望成本降低 $0.0644$ 元，使后续动作能够利用新增质量信息。

针对多工序生产，本文将 8 个零配件、3 个半成品和 1 个成品的两工序结构编码为 16 位固定策略，以有限状态期望递推精确计算交付一件合格成品的期望成本，并完全枚举 $2^{16}=65536$ 个策略。固定策略空间内的最优方案为全部零配件与半成品检测、不合格半成品拆解、成品不检测、不合格成品拆解（编码 $1111111111111101$），期望利润为 $60.2222$ 元；多方法验证结果一致。

针对问题二、问题三中次品率由抽样检测得到的情形，本文以 Beta 后验刻画参数不确定性，将样本量 $n$ 作为证据强度，并取 $n=40$ 为代表性小样本情景；以风险中性后验期望利润为准则，重新枚举问题二七变量与问题三 16 位策略。结果表明，问题二仅情形 5 补检成品，问题三最优结构保持不变、后验期望利润为 $55.03$ 元，且样本量增大时结果收敛到点估计口径。除问题二情形 5 这一策略边界情形外，不同弱信息先验下的最优策略保持一致。

\end{contestabstract}
